% =============================================================================
% Appendix D — GPIM Formalization
% Source: GPIM_Formalization_v3.md
% Companion to §1.4.1 (subsec:gpim)
% =============================================================================

\section{The Generalized Perpetual Inventory Method: Derivation and
Corporate Adjustments}
\label{app:gpim}

This appendix provides the derivations summarized in
\S\ref{subsec:gpim}. The argument proceeds in four steps: the
micro-level perpetual inventory identity, the aggregation problem under
BEA chain-weighting, the GPIM accumulation rules for current-cost and
real stocks, and the three corporate-sector adjustments that generate the
capital stock objects entering the replication.

% -------------------------------------------------------------------------
% D.1 — Micro-level identity
% -------------------------------------------------------------------------

\subsection{Micro-Level Perpetual Inventory Identity}
\label{app:gpim_micro}

For an individual asset type $i$, the standard perpetual inventory
identity defines the net capital stock recursively as
$K_{it} = IG_{it} + (1 - z_{it})\,K_{i,t-1}$,
where $IG_{it}$ is gross investment in asset $i$ and $z_{it}$ is the
asset-specific depletion rate. In current-cost terms, the inherited stock
must be marked to period-$t$ prices before the depletion deduction is
applied:
%
\begin{equation}
  \label{eq:app_micro_pim}
  K_{it}^{\mathrm{nom}}
  \;=\; p_{it}^{K}\,IG_{it}^{R}
  \;+\; (1 - z_{it})\,\frac{p_{it}^{K}}{p_{i,t-1}^{K}}\,K_{i,t-1}^{\mathrm{nom}},
\end{equation}
%
where $p_{it}^{K}$ is the price index for asset type $i$ and
$IG_{it}^{R}$ is real gross investment. Nominal depreciation for asset
$i$ is $D_{it} = z_{it}\,p_{it}^{K}\,K_{i,t-1}^{R}$.

% -------------------------------------------------------------------------
% D.2 — The aggregation problem
% -------------------------------------------------------------------------

\subsection{The Aggregation Problem Under Chain-Weighting}
\label{app:gpim_aggregation}

The BEA constructs aggregate capital stocks via Fisher-ideal
chain-aggregation of individual asset-type stocks. The standard
micro-level identity~\eqref{eq:app_micro_pim} does not survive this
aggregation. The chain index introduces a composition effect: the
aggregate price index $p_t^{K}$ reflects shifts in the asset mix, not
merely common price movements. Summing the micro identities across asset
types yields a cross-term whose magnitude depends on covariances between
individual price changes and individual stock shares --- a term that
cannot be expressed as a function of aggregate quantities alone.

This is the gap the GPIM closes. \citet[Appendix~6.5, \S V]{Shaikh2016}
shows that, to a close approximation, the aggregate current-cost stock
satisfies an accumulation rule of the same form as the micro identity,
with an aggregate composite factor replacing the asset-specific terms.
The approximation error is shown empirically to be below 0.5\% for the
US corporate sector over 1947--2005.

% -------------------------------------------------------------------------
% D.3 — GPIM accumulation rules
% -------------------------------------------------------------------------

\subsection{GPIM Accumulation Rules}
\label{app:gpim_rules}

\paragraph{Current-cost stock.}
The generalized accumulation rule for the aggregate current-cost capital
stock is:
%
\begin{equation}
  \label{eq:app_gpim_nom}
  K_t \;=\; IG_t \;+\; z_t^{\ast}\cdot K_{t-1},
  \qquad
  z_t^{\ast} \;\equiv\; (1 - z_t)\,\frac{p_t^{K}}{p_{t-1}^{K}},
\end{equation}
%
where $IG_t$ is aggregate gross investment in current prices, $z_t$ is
the theoretically correct aggregate depletion rate (defined below), and
$z_t^{\ast}$ is the composite survival-revaluation factor. The factor
compresses two operations: the fraction $(1-z_t)$ of the stock that
physically survives the period, and the revaluation of that surviving
fraction from period-$(t-1)$ to period-$t$ prices via
$p_t^{K}/p_{t-1}^{K}$.

\paragraph{Real (constant-cost) stock.}
Dividing both sides of~\eqref{eq:app_gpim_nom} by $p_t^{K}$ yields the
accumulation rule for the real stock:
%
\begin{equation}
  \label{eq:app_gpim_real}
  K_t^{R} \;=\; IG_t^{R} \;+\; (1 - z_t)\,K_{t-1}^{R},
\end{equation}
%
where $IG_t^{R} \equiv IG_t/p_t^{K}$ is real gross investment. Since
\eqref{eq:app_gpim_real} follows from~\eqref{eq:app_gpim_nom} by an
exact algebraic operation, the approximation error of the GPIM is
identical for both valuation bases.

\paragraph{Theoretically correct aggregate depletion rate.}
Aggregating the individual depreciation allowances
$D_{it} = z_{it}\,p_{it}^{K}\,K_{i,t-1}^{R}$ gives:
%
\begin{equation}
  \label{eq:app_dep_rate}
  z_t \;=\; \frac{D_t}{p_t^{K}\cdot K_{t-1}^{R}}
  \;=\; \sum_{i=1}^{n} z_{it}\,w_{it},
  \qquad
  w_{it} \;\equiv\;
  \frac{p_{it}^{K}\,K_{i,t-1}^{R}}{p_t^{K}\,K_{t-1}^{R}},
\end{equation}
%
where $w_{it}$ is the share of asset $i$'s reflated past real stock in
the aggregate. The weights use current-period asset prices $p_{it}^{K}$,
not lagged prices --- the critical distinction from the
Whelan--Liu approximation $z_t^{WL} = D_t/K_{t-1}$, which uses lagged
current-cost shares and is therefore biased by the covariance between
individual depreciation rates and asset-specific price changes.

\paragraph{Convergence properties.}
Under the constant-coefficient approximation
($z_t \approx z$, $p_t^{K}/p_{t-1}^{K} \approx 1 + g_{p_K}$), the
general solution to~\eqref{eq:app_gpim_nom} decomposes into a transient
term governed by initial conditions and a permanent term proportional to
the gross investment flow:
%
\begin{equation}
  \label{eq:app_gen_solution}
  K_t \;=\; A(z)\cdot(z^{\ast})^t \;+\; C(z)\cdot IG_t,
  \qquad
  C(z) \;\equiv\;
  \frac{1+g_I}{(1+g_I)-(1-z)(1+g_{p_K})},
\end{equation}
%
where $g_I$ is the average growth rate of gross investment and
$A(z) \equiv K_0 - C(z)\cdot IG_0$. The transient decays if and only if
$z^{\ast} < 1$, which requires:
%
\begin{equation}
  \label{eq:app_critical_z}
  z \;>\; z^{\star} \;\equiv\; \frac{g_{p_K}}{1+g_{p_K}}.
\end{equation}
%
For the US corporate sector ($g_{p_K} \approx 3.4\%$),
$z^{\star} \approx 3.3\%$. Net depreciation rates of 5--8\% place the
net stock firmly in the convergent regime; gross retirement rates below
$z^{\star}$ may leave the gross stock non-convergent over the sample
horizon. When $z > z^{\star}$, the half-life of the transient is:
%
\begin{equation}
  \label{eq:app_halflife}
  \tau_{1/2} \;=\;
  \frac{\ln 2}{\ln\!\bigl[1/\bigl((1-z)(1+g_{p_K})\bigr)\bigr]}.
\end{equation}
%
For a corporate net stock with $z \approx 0.06$, this yields
$\tau_{1/2} \approx 24$ years --- consistent with the observed narrowing
of the BEA 1993 vs.\ 2011 initial-value gap from 31\% in 1925 to
$<2\%$ by 1969.

% -------------------------------------------------------------------------
% D.4 — Shaikh's three corporate adjustments
% -------------------------------------------------------------------------

\subsection{Shaikh's Three Adjustments to BEA Official Stocks}
\label{app:gpim_adjustments}

The GPIM rules provide the algebraic machinery to generate alternative
aggregate capital stock measures by modifying BEA assumptions at three
points.

\paragraph{Adjustment 1 — Depletion rates.}
Since 1997, the BEA has assumed geometrically declining efficiency over
an infinite asset lifetime. This precludes the calculation of gross
stocks (infinitely lived assets never retire) and implies lower
depreciation rates than the pre-1997 methodology. Shaikh reverts to the
BEA 1993 finite-service-life assumption, which yields higher depreciation
rates, lower retirement rates, and time-varying $z_t$ that shifts only
with the asset mix rather than with assumed efficiency profiles.

\paragraph{Adjustment 2 — Initial values.}
The 1925 initial value of the current-cost corporate net stock is
approximately 31\% lower under BEA 1993 than under BEA 2011. By the
convergence analysis of \S\ref{app:gpim_rules}, a lower initial value
with $z > z^{\star}$ implies faster transitional growth; the gap narrows
to 10\% by 1947 and to $<2\%$ by 1969.

\paragraph{Adjustment 3 — Great Depression scrapping correction.}
The BEA methodology assumes that asset scrapping is independent of
economic conditions. Shaikh uses the IRS book-value index
\citep[Series V~115]{UScensus1975} as an external control for actual
corporate write-downs over 1929--1947. The adjustment multiplies the BEA
historical-cost stock by the IRS-to-BEA ratio over 1925--1947, then
resumes the GPIM accumulation rule from 1948 onward:
%
\begin{align}
  K_t^{\mathrm{adj}}
  &\;=\; \frac{\mathrm{IRS}_t}{\mathrm{BEA}_t}\cdot K_t^{\mathrm{BEA}},
  \qquad t \in [1925,\,1947],
  \label{eq:app_irs_adj} \\[4pt]
  K_t^{\mathrm{adj}}
  &\;=\; IG_t \;+\; z_t^{\ast}\cdot K_{t-1}^{\mathrm{adj}},
  \qquad t \geq 1948.
  \label{eq:app_gpim_resume}
\end{align}
%
The net effect: by 1947, $\texttt{KNCcorp}$ is approximately 28\% lower
than the BEA 2011 official net stock $\texttt{KNCcorpbea}$. The
adjustment is applied identically across current-cost, real, and
historical-cost valuations, since all three share the same depletion
rates.

\paragraph{Final series.}
The three adjustments plus inventory augmentation (IRS book-value ratio
extrapolated backward from 1990) produce the capital stock objects
entering the replication:

\begin{table}[!ht]
  \centering
  \small
  \caption{GPIM-adjusted capital stock series}
  \label{tab:app_gpim_series}
  \begin{tabular}{lll}
    \hline\hline
    Series & Description & Valuation \\
    \hline
    $\texttt{KGCcorp}$
      & Corporate gross fixed capital $+$ inventories (GPIM)
      & Current cost \\
    $\texttt{KNCcorp}$
      & Corporate net fixed capital (GPIM)
      & Current cost \\
    $\texttt{KNCcorpbea}$
      & Corporate net fixed capital (BEA 2011 official)
      & Current cost \\
    \hline
  \end{tabular}
  \par\vspace{0.4em}
  \footnotesize
  $\texttt{KNCcorpbea}$ is retained as a comparison benchmark.
  The ratio $\texttt{KNCcorp}/\texttt{KNCcorpbea}$ summarizes the
  cumulative effect of the three adjustments on net stock level over the
  sample. Source: \citet[Appendix~6.8]{Shaikh2016},
  \texttt{realeconomicanalysis.com}.
\end{table}

% -------------------------------------------------------------------------
% D.5 — Notation summary
% -------------------------------------------------------------------------

\subsection*{Notation}

\begin{center}
\small
\begin{tabular}{ll}
  \hline\hline
  Symbol & Definition \\
  \hline
  $K_t$ & Aggregate current-cost capital stock \\
  $K_t^{R}$ & Aggregate real (constant-cost) capital stock \\
  $IG_t$ & Aggregate gross investment (current prices) \\
  $IG_t^{R}$ & Aggregate real gross investment \\
  $z_t$ & Theoretically correct aggregate depletion rate \\
  $z_t^{\ast}$ & Composite survival-revaluation factor \\
  $z^{\star}$ & Critical depletion rate (convergence threshold) \\
  $p_t^{K}$ & Aggregate chain-weighted capital-goods price index \\
  $p_{it}^{K}$ & Asset-type $i$ price index \\
  $w_{it}$ & Reflated real-stock weight for asset $i$ \\
  $C(z)$ & Permanent-component multiplier \\
  $A(z)$ & Transient-component amplitude \\
  $\tau_{1/2}$ & Half-life of transient (initial-value) component \\
  \hline
\end{tabular}
\end{center}
